\chapter{About the written test}
\label{ch:omduggan}

The test is written and is done with pen and paper. The permitted aids are a crib sheet on one A4
sheet, both sides, and a hexadecimal calculator. \textbf{The register map is handed out with the
test} --- the indices do not have to be memorised, but reading the map and knowing what it means is
all the more necessary.

Byte sequences, checksum calculations, register packing and SPI transactions have to be doable by
hand.

\section*{Structure}

The test has two parts, of roughly equal size:

\begin{booktable}{@{}lLr@{}}
\thead{Part} & \thead{Content} & \thead{Marks} \\ \midrule
A & Your own protocol: frame design, parsing, addressing, the error model and the reliability
  chain (\kapref{ch:frames}--\kapref{ch:bussar}). & 12 \\
B & CAN and the driver: arbitration, the register map, the SPI transaction, the architecture and
  the limitations (\kapref{ch:can}--\kapref{ch:avr}). & 13 \\
\end{booktable}

A pass requires marks in \emph{both} parts. An empty part B cannot be compensated for with a
perfect part A.

\section*{Pass and distinction questions}

\textbf{The pass questions} check basic understanding: building and reading a frame, computing a
checksum, following a parser through a byte stream, reading the register map, packing a frame into
registers, decoding an SPI transaction, and explaining arbitration.

\textbf{The distinction questions} require reasoning: the reliability chain and its consequences,
analysis of faults in a byte stream, why sticky status bits are needed, the boundaries of the
layered architecture, and what the design's limitations make impossible for the driver.

\section*{The practice test}

The next chapter is a complete practice test in the same format and with the same distribution of
marks as the real one. The answers are in the chapter after that --- with the calculations, not
just the results.

Do it with pen and paper, and read the answers only afterwards.
